\par\Needspace{6\baselineskip}\medskip
\hypertarget{native-126-article093}{}\pdfbookmark[2]{93.}{article-093}
\begin{center}93.\end{center}
To make the solution of equation 15 as convenient as possible in the cases most frequent in practice, we append a special table (Table II) at the end of this work. For values of $h$ from 0 to 0,6, it gives the corresponding logarithms of $yy$, computed with the greatest care to seven decimal places. The argument $h$ increases from 0 to 0,04 in steps of one ten-thousandth, so the second differences of $\log yy$ are made to \SourceLeaf{127}vanish. Thus simple interpolation suffices in this part of the table. Since the table would have become very large had it used the same spacing throughout, from $h=0{,}04$ to the end it had to proceed in steps of one thousandth only. In this latter part we must therefore take account of the second differences if we wish to avoid errors of a few units in the seventh digit. In practice, however, smaller values of $h$ are by far the most frequent.

Both equation 15, when $h$ exceeds the limit of the table, and equation $15^{\ast}$ can be solved without difficulty by the indirect method or by other sufficiently well-known methods. It is also worth noting that a small value of $g$ can occur together with a negative value of $\cos f$ only in very eccentric orbits, as will be apparent from equation 20, to be given below in Art. 95\footnote{This equation shows that, if $\cos f$ is negative, $\varphi$ must certainly be greater than $90^0-g$.}.
